% 支持中文（自动选择字体和排版设置）
%\usepackage[UTF8]{ctex}
% 插入图片
\usepackage{graphicx}
%
\usepackage{amssymb}
\usepackage{bm}
% 颜色支持（需要这个包来使用颜色）
\usepackage{xcolor}
\usepackage{fancybox}
%
\usepackage{enumerate}
% 数学公式支持
\usepackage{amsmath}
% 超链接支持（适用于生成的PDF）
\usepackage{cite}
% \usepackage{hyperref}     % not supported by automatica
%
\usepackage[ruled,vlined,linesnumbered,norelsize,algo2e]{algorithm2e}
%
\usepackage{tabularx}
\usepackage{url}
%
% 如果需要设置参考文献的链接颜色，可以这样配置hyperref
% \hypersetup{
%     colorlinks=true,
%     linkcolor=blue,
%     citecolor=red,
%     urlcolor=magenta,
% }

% 设置行间距
\newcommand{\highlight}[1]{\textbf{\textcolor{blue}{#1}}}
\newcommand{\highlightb}[1]{%
    \colorbox{yellow}{\textbf{#1}}%
}
% 可选：带边框的版本
\newcommand{\highlightbbox}[1]{%
    \setlength{\fboxsep}{4pt}%
    \colorbox{yellow}{\parbox{\linewidth}{\color{red}\textbf{#1}}}%
    \kern-\fboxrule\fbox{\phantom{\rule{\linewidth}{0pt}}}%
}

% 可选：多种颜色选择
\newcommand{\hlred}[1]{\colorbox{red!30}{\textbf{#1}}}
\newcommand{\hlblue}[1]{\colorbox{blue!30}{\textbf{#1}}}
\newcommand{\hlgreen}[1]{\colorbox{green!30}{\textbf{#1}}}

% ---------- Boolean operators ----------
\newcommand{\LTLtrue}{\textsc{TRUE}}
\newcommand{\LTLnot}{\neg}
\newcommand{\LTLand}{\mathbin{\wedge}}
\newcommand{\LTLor}{\mathbin{\vee}}
\newcommand{\LTLimplies}{\mathbin{\rightarrow}}
\newcommand{\LTLiff}{\mathbin{\leftrightarrow}}

% Satisfaction relation (replace \models)
\newcommand{\LTLsatisfies}{\mathrel{\models}}

% Temporal operators
\newcommand{\LTLnext}{\bigcirc}
\newcommand{\LTLuntil}{\mathbin{\mathcal{U}}}
\newcommand{\LTLeventually}{\lozenge}
\newcommand{\LTLalways}{\square}

% for tabularx
\newcolumntype{L}[1]{%
>{\raggedright\arraybackslash\setlength{\parindent}{0pt}\ignorespaces}p{#1}<{\unskip}%
}

% Algorithm2e
\setlength{\algomargin}{1.5em}
%
\SetKwRepeat{Do}{do}{while}


% Theorems
% \setlength{\parindent}{1em}
% \theoremseparator{.}
%\theoremstyle{break}
% \newtheorem{problem}{Problem}%[section]
% \newtheorem{definition}{Definition}%[section]
% \newtheorem{theorem}{Theorem}%[section]
% \newtheorem{lemma}{Lemma}%[section]
% \newtheorem{proposition}{Proposition}%[section]
% %\newtheorem{proof}{Proof}[section]
% %\newenvironment{proof}{{\indent \indent \it Proof:\quad}}{\hfill $\blacksquare$\par}
% \newtheorem{remark}{Remark}%[section]
% \newtheorem{example}{Example}%[section]
% \newtheorem{assumption}{Assumption}%[section]
% \newtheorem{corollary}{Corollary}%[section]
% \newtheorem{case_study}{Case-Study}